% ===========================
% modulo_perf_timing.tex
% ===========================

\section{Módulo de Mediciones de Desempeño}% chktex 24
\label{sec:modulo-perf-timing}

\subsection{Introducción}
El \emph{Módulo de Mediciones de Desempeño} encapsula las utilidades necesarias para instrumentar el resto del sistema con mediciones temporales. Incluye un registrador buffered (CSV/JSONL), decoradores/contextos para marcar secciones críticas y utilidades de entrada/salida para analizar los archivos generados. Todas las piezas están diseñadas para ser opcionales: si la variable de entorno \texttt{PERF\_TIMINGS\_ENABLED} está desactivada, la instrumentación incurre en costo cero.

\subsection{Arquitectura de Alto Nivel}
\begin{enumerate}
    \item \textbf{Logger (\texttt{logger.py})}: escribe registros temporales en CSV (y opcionalmente JSONL) con búfer configurable.
    \item \textbf{Decoradores/Contextos (\texttt{timers.py})}: ofrece \texttt{time\_block} y \texttt{time\_section} para medir secciones de código.
    \item \textbf{Utilidades IO (\texttt{io\_utils.py})}: facilita listar, leer y filtrar archivos de timing para postprocesos o dashboards.
\end{enumerate}

\subsection{Stack tecnológico}
\begin{itemize}
    \item \textbf{Python estándar}: módulos \texttt{time}, \texttt{csv}, \texttt{json}, \texttt{pathlib}, \texttt{threading}.
    \item \textbf{Decoradores/context managers} vía \texttt{contextlib} y \texttt{functools}.
\end{itemize}

% =========================================
% DOCUMENTACIÓN POR SCRIPT
% =========================================

% ---------- logger.py ----------
\subsection{\texttt{logger.py} \- Registro buffered}
\paragraph{Descripción general}
\texttt{TimingLogger} mantiene un búfer de registros que se flushea a disco cuando alcanza un tamaño configurado o al cerrar la aplicación. Soporta formatos CSV y JSONL, creando automáticamente la carpeta \texttt{data/timings}.

\subsubsection{\texttt{class TimingRecord}}
\paragraph{Propósito} Agrupar los campos esenciales de cada medición (identificadores, marca de tiempo, duración y metadatos serializados).

\begin{listing}[H]
\caption{Conversión de un \texttt{TimingRecord} a filas CSV}
\begin{adjustbox}{max width=\textwidth}
\begin{minipage}{\textwidth}
\begin{minted}[fontsize=\small, breaklines, breakanywhere]{python}
@dataclass(frozen=True)
class TimingRecord:
    run_id: str
    epoch: int
    batch_id: int
    individual_id: int
    section: str
    start_ns: int
    end_ns: int
    duration_us: int
    extra: str

    def as_row(self) -> list[Any]:
        return [
            self.run_id,
            str(self.epoch),
            str(self.batch_id),
            str(self.individual_id),
            self.section,
            str(self.start_ns),
            str(self.end_ns),
            str(self.duration_us),
            self.extra,
        ]
\end{minted}
\end{minipage}
\end{adjustbox}
\end{listing}

\paragraph{Flujo de trabajo}
El método \texttt{TimingLogger.record(...)} normaliza los tiempos de inicio/fin, calcula la duración en microsegundos y añade el registro al búfer. Cuando el búfer alcanza \texttt{buffer\_size}, se fuerza un \emph{flush} protegido por \texttt{threading.Lock}. El cierre (vía \texttt{close}) se registra con \texttt{atexit} para seguridad.

% ---------- timers.py ----------
\subsection{\texttt{timers.py} \- Decoradores y contextos}
\paragraph{Descripción general}
Este archivo ofrece dos primitivas de instrumentación: \texttt{time\_block} para marcar bloques ad-hoc y \texttt{time\_section} para envolver funciones. Ambos respetan la bandera global \texttt{PERF\_TIMINGS\_ENABLED}, evitando costos innecesarios cuando la medición está deshabilitada.

\begin{listing}[H]
\caption{\texttt{time\_block} y \texttt{time\_section}}
\begin{adjustbox}{max width=\textwidth}
\begin{minipage}{\textwidth}
\begin{minted}[fontsize=\small, breaklines, breakanywhere]{python}
def time_block(section: str, **metadata: Any) -> _TimingBlock:
    """Context manager para cronometrar secciones concretas."""
    return _TimingBlock(section, **metadata)

def time_section(section: str, **metadata: Any) -> Callable[[F], F]:
    """Decorador que reutiliza time_block para medir funciones."""
    def decorator(func: F) -> F:
        @functools.wraps(func)
        def wrapper(*args: Any, **kwargs: Any) -> Any:
            if not PERF_TIMINGS_ENABLED:
                return func(*args, **kwargs)
            with time_block(section, **metadata):
                return func(*args, **kwargs)
        return wrapper  # type: ignore[return-value]
    return decorator
\end{minted}
\end{minipage}
\end{adjustbox}
\end{listing}

\paragraph{Resolución dinámica de metadatos}
Los decoradores aceptan \texttt{callables} para campos como \texttt{epoch} o \texttt{extra}. Las funciones auxiliares (\texttt{\_resolve\_value}, \texttt{\_to\_int}) evalúan dinámicamente estos valores usando los argumentos originales de la función envuelta.

% ---------- io_utils.py ----------
\subsection{\texttt{io\_utils.py} \- Utilidades de análisis}
\paragraph{Descripción general}
Las rutinas de \texttt{io\_utils} simplifican la inspección de resultados: listan archivos \texttt{timings\_*.csv}, validan cabeceras, convierten tipos a enteros y generan estructuras listas para consumo por analíticas o notebooks.

\paragraph{Casos de uso}
\begin{itemize}
    \item \texttt{latest\_timing\_csv()} devuelve el archivo más reciente generado por el logger.
    \item \texttt{filter\_rows(...)} permite filtrar por \texttt{run\_id}, \texttt{epoch}, \texttt{batch\_id} o \texttt{section}, facilitando dashboards de un subconjunto de métricas.
\end{itemize}

% ---------- Cobertura opcional ----------
\subsection{Cobertura y omisiones}
Para conservar el documento conciso se omiten bloques completos de \texttt{logger.py} (gestión de archivos) y \texttt{io\_utils.py} (parsers), ya que siguen patrones estándar de la biblioteca estándar. Las secciones anteriores documentan los puntos de extensión más relevantes para desarrolladores que deseen instrumentar tiempos adicionales o procesar mediciones existentes.
